\documentclass[11pt,twocolumn]{article}

% ---------- Packages ----------
\usepackage[a4paper,margin=2cm,columnsep=0.8cm]{geometry}
\usepackage{amsmath,amssymb}
\usepackage{booktabs}       % clean regression tables
\usepackage{threeparttable} % table notes
\usepackage{graphicx}
\usepackage{caption}
\usepackage{dblfloatfix}   % lets full-width tables also go at the bottom/top reliably
\captionsetup{font=small,labelfont=bf}
\usepackage{natbib}         % author-year citations
\usepackage{hyperref}
\usepackage{lipsum}        % dummy filler text: delete when writing
\usepackage{CJKutf8}
\usepackage{xurl} 

% ---------- Float placement ----------
% Let full-width tables share a page with text instead of getting a page to themselves
\renewcommand{\topfraction}{0.9}
\renewcommand{\dbltopfraction}{0.9}
\renewcommand{\textfraction}{0.07}
\renewcommand{\floatpagefraction}{0.8}
\renewcommand{\dblfloatpagefraction}{0.8}
\setcounter{topnumber}{3}
\setcounter{dbltopnumber}{3}

% ---------- Title ----------
\title{Does Air Pollution Drive Migration? An Instrumental Variables Approach with evidence from China}
\author{Sherman Goh (A0281152M) \and Bu Rui Jia, Gerianne (A0281632H)}
\date{EC3304 Research Assignment \\ Semester 1, 2026/27}

\begin{document}
\maketitle

% =====================================================
\section{Introduction}
% Motivation; research question (one sentence);
% endogeneity problem; IV approach; preview of findings.

Air pollution is a serious problem in highly industrialised countries like China. The Beijing-Tianjin-Hebei (Jing-Jin-Ji) region experienced extremely severe pollution with an annual average of fine particulate matter smaller than 2.5 micrometers ($PM_{2.5}$) of 93 micrograms per cubic meter ($\mu g/m^3$) in 2014. This value far exceeds the standard as advised by the World Health Organisation (WHO), as well as the Chinese national standards. \citep{worldbank2020china}\\

\noindent  Where health outcomes such as hospitalisation rates, pharmaceutical expenditure and mortality rates \citep{schlenker2016airports, deschenes2017defensive} are typically investigated when it comes to air pollution, its effect on migration outcomes are an interesting one. This is because re-location and migration decisions can be viewed as the internalised cost to a populace, and to better understand the willingness to pay for pollution.\citep{chen2022pollution}\\

\noindent China is a good source to investigate migration outcomes, this is partly played by its internal hukou (\begin{CJK*}{UTF8}{gbsn}户口\end{CJK*}) system, where it acts as a strict internal migration control system to divide citizens into rural and urban categories based on birthplace, and the wealth of data due to its sheer population. Because of the differences in cities, migration costs would differ, which allows for the investigation of heterogeneous effects across cities.\\

\noindent Because industrialised cities are often more polluted, with more job opportunities and a richer populace, the naive correlation between typical observed characteristics and migration outcomes are likely to be driven by endogenous  but unobserved variables, leading to a requirement for an exogenous instrument to separate its clean variation.\\

\noindent In this paper, the number of thermal inversions ($thermalnum_i$) is used as an instrument on the pollution levels ($lnpm25_i$), to investigate the effect of pollution on migration outcome. 

Something on heterogeneity...

% =====================================================
\section{Data and Context}
% Unit and timing; variable definitions; why seaport distances.
% Refer to Table~\ref{tab:summary}.
The thermal inversion IV identification strategy is carried out using cross-sectional data covering individuals' pollution exposure in 2014 and their migration decisions and demographic attributes in 2015.\\

The distances to seaports from an individual's city of origin to 3 seaports (Tianjin, Shanghai and Shenzhen) was also recorded. Because cities nearer to seaports are typically more industrialised, with more economic opportunities but also more pollution, controlling for these variables can help with removal of confounding effects on pollution and migration decisions.\\ 

These distances are geographically fixed, and are neither affected by pollution nor migration decisions. Furthermore, such distances are predetermined and not an outcome of migration decision, hence effectively exogenous.\\

% Summary statistics: fill in from Stata (e.g. summarize migrant lnpm25 thermalnum ...)
\begin{table}[!htb]
\centering
\caption{Summary Statistics ($N = 38{,}089$)}
\label{tab:summary}
\small
\begin{tabular}{lcc}
\toprule
Variable & Mean & SD \\
\midrule
migrant          & 0.102   & 0.303   \\
lnpm25           & 3.662   & 0.436   \\
thermalnum       & 199.440 & 102.394 \\
female           & 0.500   & 0.500   \\
age              & 39.932  & 11.002  \\
married          & 0.816   & 0.387   \\
urban            & 0.524   & 0.499   \\
edu\_high        & 0.155   & 0.362   \\
lntj\_dist       & 6.702   & 1.175   \\
lnsh\_dist       & 6.760   & 1.139   \\
lnsz\_dist       & 6.913   & 0.917   \\
avtemp           & 152.223 & 50.333  \\
\bottomrule
\end{tabular}
\end{table}

Table 1 provides the summary statistics for each variable found in the dataset. On first glace, we observe that the sample composes a 50-50 split between males and females, middle aged individual who are mostly married and are more unlikely to hold a college degree. This could have potential implications in the heterogenity analysis in Section 4. 

% =====================================================
\section{Empirical Strategy}
As outlined in the introduction, it is often difficult to separate the effects of pollution and the other endogenous effects such as income levels or the levels of industrialisation levels in cities on migration decisions. We first observe the regression results for the Ordinary Least Squares (OLS) Regression methods before moving on to the Instrumental Variable approach to elucidate the clean variation in pollution on migration decisions. 
\subsection{OLS Specification}
% table* spans both columns and floats to the top of the next page if needed
\begin{table*}[!t]
\centering
\begin{threeparttable}
\caption{OLS Estimates of the Effect of PM2.5 on Migration}
\label{tab:ols}
\small
\begin{tabular}{lccc}
\toprule
 & (1) & (2) & (3) \\
 & Full sample & High-skilled & Low-skilled \\
\midrule
lnpm25 & $0.0449^{***}$ & $0.0499^{***}$ & $0.0414^{***}$ \\
 & $(0.0105)$ & $(0.0158)$ & $(0.0109)$ \\
female & $-0.0105^{***}$ & $-0.0137^{*}$ & $-0.0105^{***}$ \\
 & $(0.0027)$ & $(0.0078)$ & $(0.0029)$ \\
age & $-0.0001$ & $-0.0001$ & $0.0011$ \\
 & $(0.0013)$ & $(0.0037)$ & $(0.0013)$ \\
agesq & $-0.000047^{***}$ & $-0.000056$ & $-0.000060^{***}$ \\
 & $(0.000016)$ & $(0.000048)$ & $(0.000016)$ \\
married & $0.0459^{***}$ & $0.0498^{***}$ & $0.0477^{***}$ \\
 & $(0.0051)$ & $(0.0147)$ & $(0.0052)$ \\
urban & $-0.0360^{***}$ & $-0.1128^{***}$ & $-0.0293^{***}$ \\
 & $(0.0058)$ & $(0.0179)$ & $(0.0055)$ \\
edu\_high & $0.0081$ &  &  \\
 & $(0.0080)$ &  &  \\
lntj\_dist & $0.0193^{***}$ & $0.0190^{***}$ & $0.0191^{***}$ \\
 & $(0.0038)$ & $(0.0028)$ & $(0.0047)$ \\
lnsh\_dist & $0.0043$ & $0.0137^{***}$ & $-0.0012$ \\
 & $(0.0069)$ & $(0.0029)$ & $(0.0084)$ \\
lnsz\_dist & $-0.0139^{***}$ & $0.0050$ & $-0.0209^{***}$ \\
 & $(0.0051)$ & $(0.0041)$ & $(0.0065)$ \\
Constant & $-0.0548$ & $-0.1801$ & $0.0167$ \\
 & $(0.0923)$ & $(0.1094)$ & $(0.1091)$ \\
\midrule
Observations & 38,089 & 5,903 & 32,186 \\
Clusters (cities) & 332 & 330 & 332 \\
$R^2$ & 0.031 & 0.045 & 0.032 \\
\bottomrule
\end{tabular}
\begin{tablenotes}\footnotesize
\item \textit{Notes:} Linear probability model; dependent variable is \texttt{migrant}. Standard errors clustered by city of origin (\texttt{City\_ID}) in parentheses. Columns (2) and (3) split the sample by \texttt{edu\_high}, so it is omitted as a control. $^{*}p<0.10$, $^{**}p<0.05$, $^{***}p<0.01$.
\end{tablenotes}
\end{threeparttable}
\end{table*}

% Equation, define each term; linear probability model.
\begin{align}
  \text{migrant}_{ic} &= \beta_0 + \beta_1 \,\text{lnpm25}_c \notag \\
  &\quad + \mathbf{X}_i'\boldsymbol{\gamma} + \mathbf{D}_c'\boldsymbol{\delta} + \varepsilon_{ic} \label{eq:ols}
\end{align}
% OLS results (Table~\ref{tab:ols}); why OLS is biased + expected direction of bias.
Equation 1 provides the specification for the OLS regression, where the dependent variable is the migration choice made and the independent variable is the log-pollution level. $\mathbf{X}_i'$ is a vector of  individual related control variables, while $\mathbf{D}_c'$ gives a vector of city related control variables. The list of variables are given in Table~\ref{tab:ols}\\

Table~\ref{tab:ols} gives the result for the OLS regression. The results are based on a Linear Probability Model, where the coefficient on the variable lnpm25 reflects the increase in the probability of an individual migrating out of their home city per one percent increase in the pollution level in their home city.

Cities that are more polluted are often the most industralised and could have a richer populace. (cite) A richer populace would, however, not migrate out of the city to continue to reap economic benefits gained from working in a richer city, where salaries and benefits are generally better(cite). We thus expect that this unobserved characteristic would bias the OLS result. To be specific, it is likely that the correlation between a rich city and the level of pollution is positive, while the correlation between a wealthy city and migration decision is likely to be negative. As a result the OLS estimate is likely to suffer from negative bias, which would understate the effect of pollution on migration decisions.\\

Therefore there is a need to exploit the IV approach to isolate the clean variation in pollution and to understand how it affects the migration decisions of individuals.

\subsection{Instrumental Variable Approach}
For the IV research approach, we will make use of an instrument to isolate the clean variation of pollution levels and its effect on migration decisions. The instrument used in this paper is the number of thermal inversions ($\text{thermalnum}_c$) that occured in an individual's city of origin in 2014.\\

Thermal inversions are meteorological phenomena in which the above-ground temperature is abnormally higher than the ground-level temperature(cite). Typically, when the ground level temperature is higher than the above-ground temperature, air pollutants are allowed to circulate vertically(cite). However, when this is inverted, pollutants cannot escape the atmosphere and are trapped near the ground (cite).

Therefore, under thermal inversions, city is likely to experience a higher level of pollution. 
\subsection{Identifying Assumptions}
The number of thermal inversions in a city is not a new instrument. \citet{chen2024mental} and \citet{fu2021productivity} have used thermal inversions to analyse the effects of air pollution on mental health and manufacturing firm productivity respectively. An instrument's validity hinges on 2 key facts: (1) its relevance to the dependent variable, and (2) whether it can be considered to be exogenous. 


\paragraph{Relevance.}
By the mechanics of how thermal inversions work, we can be convinced that when it happens, the city is likely to experience a higher effective level of pollution. As such we can expect that the higher number of thermal inversions in an individual's origin city, it is likely that the individual also experiences a higher level of pollution. We can formalise this to say that the number of thermal inversions are positively correlated to the pollution level in the city and thus $Cov(\text{thermalnum}_c,\ln{pm25})>0$.
\begin{align}
  \text{lnpm25}_c &= \pi_0 + \pi_1 \,\text{thermalnum}_c \notag \\
  &\quad + \mathbf{X}_i'\boldsymbol{\lambda} + \mathbf{D}_c'\boldsymbol{\theta} + u_{ic} \label{eq:fs}
\end{align}
This is further supported by the evidence in the data. Equation \ref{eq:fs} gives the first stage regression of the pollution level on the instrument, with the same controls found in Table~\ref{tab:ols}. When regressed on the entire dataset, the First stage F statistic (cluster robust) is 52.13. Similarly, when the dataset was split into high and low skilled workers the first stage F statistic (cluster robust) gives 32.42 and 51.41 respectively.\\

Not only is the instrument relevant, it is also strong, with F statistics well over 10. With a strong first stage, the IV results could be highly convincing.

\begin{table*}[!t]
\centering
\begin{threeparttable}
\caption{2SLS Estimates and First Stage: Thermal Inversions as an Instrument for PM2.5}
\label{tab:iv}
\small
\begin{tabular}{lccc}
\toprule
 & (1) & (2) & (3) \\
 & Full sample & High-skilled & Low-skilled \\
\midrule
\multicolumn{4}{l}{\textit{Panel A: Second stage (dependent variable: migrant)}} \\
lnpm25 & $0.1108^{***}$ & $0.1676^{***}$ & $0.1030^{***}$ \\
 & $(0.0292)$ & $(0.0507)$ & $(0.0308)$ \\
\midrule
\multicolumn{4}{l}{\textit{Panel B: First stage (dependent variable: lnpm25)}} \\
thermalnum & $0.00183^{***}$ & $0.00151^{***}$ & $0.00186^{***}$ \\
 & $(0.00025)$ & $(0.00026)$ & $(0.00026)$ \\
First-stage $F$ (cluster-robust) & 52.13 & 32.42 & 51.42 \\
Partial $R^2$ & 0.198 & 0.145 & 0.201 \\
\midrule
Observations & 38,089 & 5,903 & 32,186 \\
\bottomrule
\end{tabular}
\begin{tablenotes}\footnotesize
\item \textit{Notes:} \texttt{lnpm25} is instrumented with \texttt{thermalnum}. All regressions control for \texttt{female}, \texttt{age}, \texttt{agesq}, \texttt{married}, \texttt{urban}, \texttt{lntj\_dist}, \texttt{lnsh\_dist}, \texttt{lnsz\_dist}, and \texttt{edu\_high} in column (1). Standard errors clustered by city of origin in parentheses. The first-stage $F$ is the cluster-robust $F$ on the excluded instrument. $^{*}p<0.10$, $^{**}p<0.05$, $^{***}p<0.01$.
\end{tablenotes}
\end{threeparttable}
\end{table*}

\paragraph{Exogeneity}
Thermal inversions are of meteorological origin, it is therefore unpredictable and essentially random. The presence of an inversion is therefore not at all correlated to economic activity nor how wealthy a country is. \\

Thermal inversions affecting migration decisions through other channels would be a threat to exogeneity. \citet{arceo2016} found that thermal inversions themselves do not pose a health risk, other than the accumulation of pollutants. Therefore, in respect of this channel of human health deterioration, it is unlikely that thermal inversions causes a human health to deteriorate forcing them to migrate. It is thus plausible to believe that thermal inversions do ont have a direct channel other than through pollution to affect migration decisions, and we hence establish the validity of this instrument. 

However, the thermal inversion IV is constructed from ground-level and above-ground temperatures, in this vein, ground level temperatures may hence affect migration decisions directly. This problem in the instrument will be addressed in Section 4.2

% Temperature channel; conditional assumption from part (c).

% =====================================================

% Full-width tables appear at the top of the page AFTER the one where this code sits.
% Code is placed one page early so each table lands on the page where its discussion starts.
% Once your real text is in, move these blocks up/down if a table lands too early or late.

\section{Results}
\begin{table*}[!t]
\centering
\begin{threeparttable}
\caption{2SLS Estimates Controlling for Average Ground-Level Temperature}
\label{tab:iv_temp}
\small
\begin{tabular}{lccc}
\toprule
 & (1) & (2) & (3) \\
 & Full sample & High-skilled & Low-skilled \\
\midrule
\multicolumn{4}{l}{\textit{Panel A: Second stage (dependent variable: migrant)}} \\
lnpm25 & $0.1158^{***}$ & $0.1627^{***}$ & $0.1059^{***}$ \\
 & $(0.0262)$ & $(0.0443)$ & $(0.0275)$ \\
avtemp & $0.00013$ & $-0.00009$ & $0.00008$ \\
 & $(0.00015)$ & $(0.00020)$ & $(0.00018)$ \\
\midrule
\multicolumn{4}{l}{\textit{Panel B: First stage (dependent variable: lnpm25)}} \\
thermalnum & $0.00209^{***}$ & $0.00176^{***}$ & $0.00211^{***}$ \\
 & $(0.00023)$ & $(0.00024)$ & $(0.00023)$ \\
avtemp & $0.0037^{***}$ & $0.0032^{***}$ & $0.0039^{***}$ \\
 & $(0.0007)$ & $(0.0007)$ & $(0.0007)$ \\
First-stage $F$ (cluster-robust) & 86.18 & 52.58 & 85.32 \\
Partial $R^2$ & 0.261 & 0.203 & 0.265 \\
\midrule
Observations & 38,089 & 5,903 & 32,186 \\
\bottomrule
\end{tabular}
\begin{tablenotes}\footnotesize
\item \textit{Notes:} As in Table~\ref{tab:iv}, with \texttt{avtemp} added as a control in both stages. Standard errors clustered by city of origin in parentheses. $^{*}p<0.10$, $^{**}p<0.05$, $^{***}p<0.01$.
\end{tablenotes}
\end{threeparttable}
\end{table*}

\begin{table*}[!t]
\centering
\begin{threeparttable}
\caption{Heterogeneity in the Effect of PM2.5 on Migration (2SLS)}
\label{tab:het}
\small
\setlength{\tabcolsep}{4pt}
\begin{tabular}{lcccccc}
\toprule
 & \multicolumn{2}{c}{Age} & \multicolumn{2}{c}{Marital status} & \multicolumn{2}{c}{Birthplace} \\
\cmidrule(lr){2-3}\cmidrule(lr){4-5}\cmidrule(lr){6-7}
 & (1) & (2) & (3) & (4) & (5) & (6) \\
 & $>$35 & $\le$35 & Married & Unmarried & Urban & Rural \\
\midrule
\multicolumn{7}{l}{\textit{Panel A: Second stage (dependent variable: migrant)}} \\
lnpm25 & $0.0804^{***}$ & $0.1674^{***}$ & $0.1159^{***}$ & $0.1137^{***}$ & $0.0931^{***}$ & $0.1332^{***}$ \\
 & $(0.0236)$ & $(0.0346)$ & $(0.0289)$ & $(0.0252)$ & $(0.0231)$ & $(0.0357)$ \\
\midrule
\multicolumn{7}{l}{\textit{Panel B: First stage (dependent variable: lnpm25)}} \\
thermalnum & $0.00200^{***}$ & $0.00224^{***}$ & $0.00205^{***}$ & $0.00223^{***}$ & $0.00186^{***}$ & $0.00223^{***}$ \\
 & $(0.00022)$ & $(0.00024)$ & $(0.00021)$ & $(0.00030)$ & $(0.00022)$ & $(0.00026)$ \\
First-stage $F$ (cluster-robust) & 82.69 & 87.86 & 92.44 & 53.97 & 73.91 & 74.35 \\
\midrule
Observations & 23,763 & 14,326 & 31,098 & 6,991 & 19,976 & 18,113 \\
\bottomrule
\end{tabular}
\begin{tablenotes}\footnotesize
\item \textit{Notes:} All columns are 2SLS with \texttt{lnpm25} instrumented by \texttt{thermalnum} and include the controls from Table~\ref{tab:iv_temp}, including \texttt{avtemp}. Columns (1)--(2) keep \texttt{age} and \texttt{agesq}; columns (3)--(6) drop \texttt{married}, since it is constant within each subsample. Columns (5) and (6) are restricted to married individuals and split by \texttt{urban}, which is also dropped. Standard errors clustered by city of origin in parentheses. $^{*}p<0.10$, $^{**}p<0.05$, $^{***}p<0.01$.
\end{tablenotes}
\end{threeparttable}
\end{table*}

To do this, we use a 2 Stage Least Squares preocedure. Equation 2 gives the first stage where the dependent variable is regressed on the instrument $\text{thermalnum}_c$

\begin{align}
  \text{migrant}_{ic} &= \beta_0 + \beta_1 \,\widehat{\text{lnpm25}}_c \notag \\
  &\quad + \mathbf{X}_i'\boldsymbol{\gamma} + \mathbf{D}_c'\boldsymbol{\delta} + \varepsilon_{ic} \label{eq:ss}
\end{align}

\subsection{IV Results}
% Full-width tables appear at the top of the page AFTER the one where this code sits.
% Code is placed one page early so each table lands on the page where its discussion starts.
% Once your real text is in, move these blocks up/down if a table lands too early or late.



% First stage + IV, Table~\ref{tab:iv}; compare with OLS.

\lipsum[1]

\subsection{Controlling for Ground Temperature}
We alluded earlier to the fact that ground level conditions could affect the thermal inversion instrument. The average ground level temperature is used to construct the thermal inversion instrument, and therefore is directly correlated to the instrument. Not controlling for this variable will result in this variable residing the in the error term in both the first stage and second stage regressions, resulting in omitted variable bias. Therefore, we propose that the first stage be modified as in Equation 4.
\begin{align}
  \text{lnpm25}_c &= \pi_0 + \pi_1 \,\text{thermalnum}_c + \pi_2 \,\text{avtemp}_c \notag \\
  &\quad + \mathbf{X}_i'\boldsymbol{\lambda} + \mathbf{D}_c'\boldsymbol{\theta} + u_{ic}
\end{align}

Compared to Equation 2, this first stage equation includes $\text{avtemp}_c$ as an additional control variable. The new conditional assumption for exogeneity to hold in this case is formally $Cov(\text{thermalnum}_c,u\ |\  \mathbf{X}_i',\ \text{avtemp}_c )=0$.\\

The regression results in Table~\ref{tab:iv_temp} provide good evidence for the need to control for the average ground temperature. Firstly, as compared to the results in Table~\ref{tab:iv}, the coefficient on lnpm25 for the full sample remains significant at the 1\% level and barely moved (0.1108 in Table 3 and 0.1158 in Table 4). Also, once average ground temperature has been controlled for, there is no evidence of a ground level temperature effect on migration decisions, once pollution has been accounted for. This is observed through the statitically insignificant coefficent on avtemp in the second stage regression across all three regressions in Table 4.

Next, in the first stage regression in Table~\ref{tab:iv_temp}, we see that the coefficient on avtemp is statistically significant (0.0037) at the 1\% level across all three regressions. This suggests that indeed, the ground level temperature is indeed correlated with the level of pollution. Finally, we also observe that the F statistic increases from about 30-50 in Table 3 to about 50-80 across all 3 regressions in Table 4. This shows that controlling for the average temperature improves the relevance and the strength of the instrument. 

\subsection{Heterogeneity}


% Part (d): age, marital status, urban/rural among married.
% Table~\ref{tab:het}.

\lipsum[1]

% =====================================================
\section{Policy Implications}

\lipsum[1]

% =====================================================
\section{Limitations and Conclusion}

\lipsum[1]

% =====================================================
\bibliographystyle{apalike}
\bibliography{references} % create references.bib in Overleaf



\end{document}
